

\section{Preliminaries}


\subsection{Notations}


We denote an unweighted, undirected graph as $G = (V, E)$, where $V = \{v_1, v_2, \ldots, v_N\}$ is the set of nodes, and $E \subseteq V \times V$ is the set of edges.
Self-loops, i.e., edges of the form $(u,u)$, are not allowed.
Let $N = |V|$ and $M = |E|$ denote the number of nodes and edges in $G$, respectively. 
For a node $v_i \in V$, 
let $\mathcal{E}(i)$ be the set of incident edges of node $v_i$, and
let $\mathcal{N}(i)$ denote the indices of its neighbors, where $\mathcal{N}(i) = \{j \mid (v_i, v_j) \in \mathcal{E}(i)\}$. 
The degree of $v_i$ is defined as $d(v_{i}) = |\mathcal{N}(i)|$, and $d(G)$ denotes the maximum degree in $G$. 
If two graphs $G = (V, E)$ and $G' = (V', E')$ differ by exactly one edge $e = (v_i, v_j)$, 
we write $G \sim_e G'$, or simply $G \sim G'$.



\subsection{Graph Pattern Counting}
\label{subsec:prelim_pattern_counting}



Given any acyclic graph pattern, we use $k$ to denote the number of edges in the pattern. 
Let $Q: \mathcal{G} \to \mathbb{Z}_{\geq 0}$ be the pattern counting query that returns the number of subgraphs in an input graph $G$.
A common type of acyclic pattern is a path, which is a sequence of distinct nodes connected by edges. 
% In traditional graph analytics, path counting is often discussed alongside walk counting.
In contrast, walks may include repeated nodes, making them computationally easier to count in the non-private setting: walk counting is solvable in polynomial time, whereas path counting is provably \#W[1]-complete~\cite{flum2004parameterized}.
We refer to a path with $k$ edges as a $k$-line path, and a walk with $k$ edges as a $k$-line walk.
Note that walk counting is not considered part of acyclic pattern counting, as walks can contain cycles, e.g., a $3$-line walk can form a triangle.


% \red{Figure 1 here???}



% \subsection{Differential Privacy on Graphs}

% \subsection{Local Differential Privacy on Graphs}

% Local DP definition, 




% that are isomorphic to $S$, i.e., 
% \begin{equation*}
% Q_S(G) := | \{ H \subseteq G \mid H \cong S\}|,    
% \end{equation*}
% where $H \cong S$ means there exists a bijection $\phi: V_S \to V_H$ such that $(u_i,u_j) \in E_S$ if and only if $(\phi(u_i), \phi(u_j)) \in E_H$.
% We refer to $\phi(u_i)$ as the image of $u_i$ under the bijection $\phi$.

% We use $\sigma(S)$ to denote the symmetry factor of $S$, defined as the number of bijections $\phi_S: V_S \to V_S$ such that $(u,v) \in E_S$ if and only if $(\phi_S(u_i), \phi_S(u_j)) \in E_S$ for all $u_i, u_j \in V_S$.
% Intuitively, $\sigma(S)$ counts the distinct node relabelings of $V_S$ that leave $S$ unchanged; for example, if $S$ is a path, then $\sigma(S) = 2$.
% Consequently, if $H \cong S$, there are exactly $\sigma(S)$ bijections $\phi: V_S \to V_H$ satisfying the isomorphism condition.



% O and regarding both \eps and \beta as a constant 


% $\{V_1, V_2, \cdots, V_k\}$ is a partition of $V$ if, $V_i \subseteq V$, $\cup V_i = V$ and any two node subsets are disjoint. 
% We use $E(V_i, V_j)$ to 

% node set partition




\subsection{Differential Privacy on Graphs}

% Most DP works on graph analytics study the central setting, where a data curator has access to the whole ... He is responsible to ... to make the graph analytical results do not leak some sensitive information.  

% Two graphs are ... node-neighbors 

% Edge neighbors.

% A mechanism is called to be DP mechanism if...

% BY using ... or ..., we have definitions for node-DP or edge-DP.

% DP satisfies several deserived properties...



Most works on DP in graph analytics operate under the central-DP setting, where a trusted data curator has access to the entire graph dataset. 
Upon receiving a query, the curator performs analytics and ensures that the output satisfies DP. 
This guarantees that one cannot distinguish between two neighboring graphs based on the output, thereby preserving privacy.
Formally, a randomized mechanism is said to satisfy DP if the following holds:

\begin{definition}[Differential Privacy]
\label{def:dp}
A randomized mechanism $\mathcal{M}: \mathcal{G} \to \mathcal{K}$ satisfies $\varepsilon$-DP if, for any pair of neighboring graphs $G, G'$ and any measurable subset $\mathcal{S} \subseteq \mathcal{K}$, it holds that
\begin{equation*}
    \Pr[\mathcal{M}(G) \in \mathcal{S}] \leq e^{\varepsilon} \Pr[\mathcal{M}(G') \in \mathcal{S}].
\end{equation*}
\end{definition}


Two graphs $G$ and $G'$ are said to be \textit{node-neighboring} if they differ by exactly one node and all of its incident edges. 
Similarly, they are \textit{edge-neighboring} if they differ by exactly one edge. 
Applying these notions to the DP definition yields the concepts of \textit{node-differential privacy} (node-DP) and \textit{edge-differential privacy} (edge-DP)~\cite{hay2009accurate}, respectively. 

DP satisfies several desirable properties, as outlined below.


\begin{lemma}[Post-Processing]
\label{lem:pp}
Let $\mathcal{M}: \mathcal{G} \to \mathcal{K}$ be a randomized mechanism that satisfies $\varepsilon$-DP. 
Then for any (possibly randomized) function $\mathcal{F}: \mathcal{K} \to \mathcal{Z}$, the mechanism $\mathcal{F} \circ \mathcal{M}: \mathcal{G} \to \mathcal{Z}$ also satisfies $\varepsilon$-DP. 
\end{lemma}

\begin{lemma}[Basic Composition] 
\label{lem:bc}
Let $\mathcal{M}_1, \mathcal{M}_2, \dots, \mathcal{M}_n$ be a series of randomized mechanisms, 
where for each $i \in \{1, 2, \dots, n\}$,
$\mathcal{M}_i: \mathcal{G} \to \mathcal{K}_i$ satisfies $\varepsilon$-DP. 
The composition mechanism
$
\mathcal{M}(G) = (\mathcal{M}_1(G), \dots, \mathcal{M}_n(G))
$
satisfies $n\varepsilon$-DP.
\end{lemma}

\begin{lemma}[Parallel Composition]
\label{lem:pc}
Let $\mathcal{M}_1, \mathcal{M}_2, \ldots, \mathcal{M}_n$ be mechanisms, each satisfying $\varepsilon$-DP. 
If for any pair of neighboring graphs $G$ and $G'$, there exists at most one mechanism $\mathcal{M}_i$ such that $\mathcal{M}_i(G) \neq \mathcal{M}_i(G')$, then the composition mechanism
$
\mathcal{M}(G) = (\mathcal{M}_1(G), \dots, \mathcal{M}_n(G))
$
satisfies $\varepsilon$-DP.
\end{lemma}



A common DP mechanism is the \textit{Laplace mechanism}. 
Here, for any arbitrary graph query $Q: \mathcal{G} \to \mathcal{R}$, its \textit{global sensitivity} is defined as
\begin{equation*}
    \mathrm{GS}_Q := \max_{G, G'} \| Q(G) - Q(G') \|_1,
\end{equation*}
where $G$ and $G'$ are any pair of neighboring graphs. Then, the Laplace mechanism is defined as follows:
\begin{definition}[Laplace Mechanism~\cite{dwork2006calibrating}]
\label{def:lap_mech}
Given a graph query $Q: \mathcal{G} \to \mathcal{R}$, the mechanism
\begin{equation*}
    \mathcal{M}(G) = Q(G) + \operatorname{Lap}\left(\frac{1}{\varepsilon}\mathrm{GS}_Q\right)
\end{equation*}
satisfies $\varepsilon$-DP, where $\operatorname{Lap}(\mathrm{GS}_Q/\varepsilon)$ denotes a Laplace random variable with mean $0$ and scale $\mathrm{GS}_Q/\varepsilon$.
\end{definition}
Specifically, computing global sensitivity under node-neighboring or edge-neighboring graphs leads to node-DP or edge-DP guarantees, respectively.





\subsection{Local Differential Privacy}
\label{subsec:ldp}



In contrast to the central-DP setting, where data analytics is performed centrally by a trusted curator, 
LDP studies a distributed setting,
in which each party corresponds to a node with access only to its own information. 
When performing data analytics tasks, each node first privatizes its data locally, then sends the noisy output to an analyzer for post-processing.
To guarantee LDP, we must ensure that the collection of messages received by the analyzer satisfies DP, as defined in Definition~\ref{def:dp}.
More precisely, each node $v_i$ has access to a subgraph $G_i = \left(\{v_i\} \cup \{v_j \mid j \in \mathcal{N}(i)\}, \mathcal{E}(i)\right)$.
To perform graph analytics,
each $v_i$ invokes 
a local randomizer $\mathcal{R}$ on $G_i$, which produces a privatized output $\mathcal{X}_i = \mathcal{R}(G_i)$.
The analyzer $\mathcal{A}$ then aggregates the local outputs $(\mathcal{X}_1, \mathcal{X}_2, \dots, \mathcal{X}_N)$ to produce the final result.
A mechanism is said to satisfy $\varepsilon$-LDP if the joint output $(\mathcal{X}_1, \mathcal{X}_2, \dots, \mathcal{X}_N)$ satisfies $\varepsilon$-DP.
Similar to the central-DP setting, applying node- or edge-neighboring definitions yields notions of \textit{node-LDP} and \textit{edge-LDP}, respectively.
To date, all existing LDP graph analytics literature~\cite{sun2019analyzing,ye2020towards, ye2020lf,Imola2021, Imola2022,eden2023triangle, he2025robust} has focused on the edge-LDP setting, which is also the focus of our work.




A commonly used mechanism under LDP is RR, in which each node $v_i \in V$ outputs a randomized neighbor list. 
The analyzer $\mathcal{A}$ can then construct a noisy graph based on the received neighboring lists and perform graph analytics on it.
Since each edge is shared between two incident nodes, to avoid duplication, the information of edge $(v_i,v_j)$ is reported only by node $v_i$ when $i < j$.
Specifically, for each node $v_i$, RR takes as input a binary neighbor list $\textbf{a}_i = (a_{i,1}, \dots, a_{i,i-1}) \in \{0, 1\}^{i-1}$, where $a_{i,j}$ is $1$ if $j \in \mathcal{N}(i)$ and $0$ otherwise. 
For each entry $a_{i,j}$, RR keeps the true bit with probability $p ={e^\varepsilon}/({e^\varepsilon + 1})$ and flips it with probability $q = {1}/({e^\varepsilon + 1})$. 
This mechanism satisfies $\varepsilon$-edge-LDP.






\paragraph{Multi-round LDP mechanism}
% \red{The above model works for some simple queries like edge counting~\cite{} and star-counting~\cite{???}, where we only require... However, for more complex graph analytical tasks,... we require... Here, we have $r$-round communications. In $\ell$-th round with $\ell \in\{1,\dots, r\}$, each node $v_i$ invoke $\mathcal{R}^{(\ell)}$..., where $\mathcal{R}^{(\ell)}$ can utilize the results of all first $\ell-1$ rounds of all parties. To guarantee the while process preserve $\varepsilon$..., for any $\ell$ we ensure ... preserves ..., then by applying composition, we have the whole processes preserves ... }


The above basic LDP model is sufficient for simple queries such as edge counting and star counting~\cite{Imola2021}, where the final result can be obtained in a single round. 
However, more complex graph analytics tasks like triangle counting~\cite{Imola2021, Imola2022,he2025robust} often require collaboration between nodes and the analyzer across multiple rounds of interaction.
Here, we consider a $k$-round LDP mechanism.
In round $\ell \in \{1, \dots, k\}$, each node $v_i$ invokes a round-specific randomizer $\mathcal{R}^{(\ell)}$ on its local subgraph $G_i$, possibly incorporating messages released by the other parties in the previous $\ell-1$ rounds.
Let $\mathcal{X}^{(\ell)}_i$ denote the privatized output of node $v_i$ in round $\ell$, and let the collection of all node messages in round $\ell$ be
$\mathcal{X}^{(\ell)} = 
    (\mathcal{X}^{(\ell)}_1, \ldots, \mathcal{X}^{(\ell)}_N)$.
The mechanism satisfies $\varepsilon$-LDP if the joint messages from all rounds, i.e., $(\mathcal{X}^{(1)}, \dots, \mathcal{X}^{(k)})$, satisfy $\varepsilon$-DP.
To ensure this, we can allocate the privacy budget across rounds such that each $\mathcal{X}^{(\ell)}$ satisfies $\varepsilon / k$-DP. 
By applying basic composition (Lemma~\ref{lem:bc}), the entire $k$-round process satisfies $\varepsilon$-DP as required.
To quantify the communication cost incurred across rounds, we sum the sizes (in bits) of all messages transmitted. 
For computational cost, we report the total cost incurred across all rounds by the analyzer and by a single node.





\subsection{Useful Concentration Bounds}
The following three lemmas provide high-probability concentration bounds that will be used in the subsequent analysis.
 
\begin{lemma}[Concentration Bound of Laplace Distributions~\cite{chan11continual}]
\label{lem:con_lap}
Let $\eta_1, \eta_2, \dots, \eta_k$ be independent random variables with $\eta_i \sim \mathrm{Lap}(b_i)$. Then, for any $\beta > 0$,
\begin{equation*}
\label{eq:lap_bound}
\Pr\left[\bigg| \sum_i \eta_i \bigg| \geq \sqrt{8 \sum_i b_i^2 \cdot \log\frac{2}{\beta}}\right] \leq \beta.
\end{equation*}

\end{lemma}

 
\begin{lemma}[Multiplicative Chernoff Bound~\cite{doerr2019probabilistic}]
\label{lem:con_cher}
Let \(X_1,\dots,X_n\) be independent random variables taking values in \([0,1]\).
Let \(X=\sum_{i=1}^n X_i\).
For any \(0 < \delta \le 1\),
\begin{equation*}
    \Pr[X \ge (1+\delta)\,\mathbb{E}[X]] 
  \leq \exp(-\frac{1}{3}\delta^{2}\,\mathbb{E}[X]).
\end{equation*}
    
\end{lemma}

\begin{lemma}[Chebyshev’s inequality~\cite{chebyshev1867valeurs}]
\label{lem:che_neq}
Let $X$ be a random variable with $\mathbb{E}[X]=\mu$ and $\operatorname{Var}(X)=\sigma^{2}$.  
For any $k>0$,
\begin{equation*}
  \Pr[|X-\mu|\ge k\sigma]\le\frac{1}{k^{2}} .
\end{equation*}
\end{lemma}


 


